% META: {"id":"1SPE-EXPONENTIELLE-RE-C5","chapitre":"1SPE-EXPONENTIELLE","type_objet":"remediation","capacites_codes":["C5"],"status":"approved","origin":"generated","status_history":["generated","audited","accepted","approved"],"release_acceptance":"RELEASE_OWNER_BATCH_ACCEPTANCE","acceptance_closure_digest":"sha256:9b3ccf9a81c5520fbb7e03b7b2d3e7bf2057a02834b6d908bf3be5f49f3b553f","acceptance_receipt_digest":"sha256:4fad86f0282e0fa4e0419cca1ad6379ae7af5a280c04a4a90880f90a1c044242"}

\subsection*{Remédiation C5 — Lire et exploiter un modèle exponentiel}

\begin{exercice}{1SPE-EXPO-RE-C5-EX1}{1}{7}
Une quantité est modélisée par $Q(t)=800\mathrm{e}^{0{,}03t}$ pour $t\geqslant0$.
\begin{enumerate}
\item Calculer $Q(0)$ et $Q(10)$, arrondi à l'unité.
\item Calculer $Q'(t)$ et interpréter son signe.
\end{enumerate}
\end{exercice}
\begin{corrige}{1SPE-EXPO-RE-C5-EX1}
$Q(0)=800$ et $Q(10)=800\mathrm{e}^{0{,}3}\approx1080$. De plus, $Q'(t)=24\mathrm{e}^{0{,}03t}>0$ : le modèle est croissant.
\end{corrige}

\begin{exercice}{1SPE-EXPO-RE-C5-EX2}{2}{8}
Une concentration est modélisée par $C(t)=20\mathrm{e}^{-0{,}2t}$.
\begin{enumerate}
\item Calculer $C(0)$, $C(5)$ et $C(10)$ à $10^{-2}$ près.
\item Montrer que $C(t+5)=\mathrm{e}^{-1}C(t)$ et interpréter.
\item Choisir, parmi trois courbes fournies par le professeur, celle qui représente ce modèle et justifier.
\end{enumerate}
\end{exercice}
\begin{corrige}{1SPE-EXPO-RE-C5-EX2}
$C(0)=20$, $C(5)=20\mathrm{e}^{-1}\approx7{,}36$ et $C(10)=20\mathrm{e}^{-2}\approx2{,}71$. La relation demandée traduit un facteur multiplicatif constant $\mathrm{e}^{-1}$ toutes les cinq unités de temps. La courbe correcte est positive, passe par $(0;20)$ et décroît.
\end{corrige}

% BEGIN-VERIFY
% from math import exp
% assert round(800*exp(0.3)) == 1080
% assert abs(20*exp(-1)-7.36) < 0.01
% assert abs(20*exp(-2)-2.71) < 0.01
% END-VERIFY
